\begin{proposition}[A lower bound on the faithful action image]
\label{B:four-by-four-image24}
If a finite monoid has a $4\times4$ rectangular-band minimal
ideal and classifying space homotopy equivalent to $\mathbb CP^2$,
its image in the monoid of row and column transformation pairs
has at least twenty-four elements.  This concerns the cardinality
of the image, with no upper bound on the source monoid.
\end{proposition}
\begin{proof}
Proposition~\ref{B:four-by-four-CP2-hulls} reduces the image to
one of the three allowed hulls $217,216,221$.  The image fixes
the band, and therefore has the same integral second homology
as the source.  Images of order at most twenty-two have zero
pulled-back rational square by the explicit fourth-homology and
cup certificates in Proposition~\ref{B:four-by-four-bounded22}.

For order twenty-three, run the same closed-subset enumeration
with six optional actors instead of five.  No search rule changes:
every intermediate closure is contained in the final image, so
the bounded closure procedure remains exhaustive.  The exact
counts are
\[
 \begin{array}{c|r|r|r|r}
 \text{hull}&\text{rank-one images}&\text{permutation/opposite orbits}
 &H_4(\mathbb F_2)\ne0&H_3=0,\ H_4\ne0\\\hline
 217&65464&4135&294&3\\
 216&5975&2992&188&0\\
 221&14522&3661&287&1.
 \end{array}
\]
The last two columns concern the orbit representatives and use
$\mathbb F_2$ coefficients.  The literal input tables and
complete lists are the files with prefixes
\path{code/outputs/d_cp23_hull_217},
\path{code/outputs/d_cp23_hull_216}, and
\path{code/outputs/d_cp23_hull_221}.
The binary bar calculation checks every consecutive differential
composite before taking ranks.

Zero fourth binary homology makes every rational square zero.
Nonzero third binary homology is now an obstruction even for
a nonfaithful source: Lemma~\ref{B:band-quotient-H3-surjection}
makes the source third homology surject onto that of its image.
Thus only the four representatives in the last column remain.
Their identifiers and exact integral projective-prefix ranks are
\[
 \begin{array}{c|c|ccccc}
 \text{hull}&\text{case}&P_0&P_1&P_2&P_3&P_4\\\hline
 217&1539&4&14&37&76&150\\
 217&1540&4&14&37&76&150\\
 217&4019&4&14&32&66&140\\
 221&3653&4&12&17&26&60.
 \end{array}
\]
Their coinvariant prefixes have dimensions $(1,0,1,0,2)$.
The integral affine hull cocycle pairs with the second-degree
generator by $-1,-1,1,-1$, respectively.  Its cup square is
identically zero on every degree-four projective basis vector
under the actual integral comparison to the normalized bar.
This proves square-zero cohomology, independently of any
uncomputed resolution tail.

The producer \path{code/a_cp23_image_cups.py} supplies the full
prefixes and sparse comparisons.  The standalone verifier
\path{code/a_verify_cp23_image_cups.py} imports no resolution
or comparison producer.  It reconstructs every projective,
checks integral exactness and saturation through $P_3$, all
equivariance and chain equations, the primitive degree-two
pairing, and the zero square on the full $P_4$ basis.  The
independent results are recorded in
\path{code/outputs/a_cp23_image_cups_independently_verified.json}.
Every possible image of order at most twenty-three is therefore
excluded, proving the proposition.
\end{proof}
